\documentclass[11pt,a4paper]{article}

\usepackage[margin=1in]{geometry}
\usepackage{amsmath,amssymb,amsthm,mathtools}
\usepackage{microtype}
\emergencystretch=3em
\usepackage[T1]{fontenc}
\usepackage[utf8]{inputenc}
\usepackage{lmodern}
\usepackage{booktabs}
\usepackage{enumitem}
\usepackage[dvipsnames]{xcolor}
\usepackage{hyperref}
\usepackage[nameinlink,capitalize]{cleveref}

\hypersetup{
  colorlinks=true,
  linkcolor=blue!50!black,
  citecolor=blue!50!black,
  urlcolor=blue!50!black,
  pdftitle={Proofs of Fourier-optics envelope identities in aberration-corrected LEEM},
  pdfauthor={Wilson Lau}
}

\newtheorem{theorem}{Theorem}[section]
\newtheorem{lemma}[theorem]{Lemma}
\newtheorem{proposition}[theorem]{Proposition}
\newtheorem{corollary}[theorem]{Corollary}
\theoremstyle{definition}
\newtheorem{definition}[theorem]{Definition}
\newtheorem{approximation}[theorem]{Approximation}
\theoremstyle{remark}
\newtheorem{remark}[theorem]{Remark}

\DeclareMathOperator{\ReOp}{Re}
\DeclareMathOperator{\ImOp}{Im}
\DeclareMathOperator{\Arg}{Arg}
\DeclareMathOperator{\diag}{diag}

\newcommand{\R}{\mathbb{R}}
\newcommand{\C}{\mathbb{C}}
\newcommand{\Z}{\mathbb{Z}}
\newcommand{\dd}{\mathrm{d}}
\newcommand{\ee}{\mathrm{e}}
\newcommand{\ii}{\mathrm{i}}
\newcommand{\conj}[1]{\overline{#1}}
\newcommand{\abs}[1]{\lvert#1\rvert}
\newcommand{\norm}[1]{\lVert#1\rVert}
\newcommand{\qap}{q_{\mathrm{ap}}}
\newcommand{\qill}{q_{\mathrm{ill}}}
\newcommand{\sill}{\sigma_{\mathrm{ill}}}
\newcommand{\sE}{\sigma_E}
\newcommand{\DE}{\Delta E}
\newcommand{\Dz}{\Delta z}
\newcommand{\EStot}{E_{S}}
\newcommand{\ECtot}{E_{C,\mathrm{tot}}}
\newcommand{\ECC}{E_{CC}}
\newcommand{\RFO}{R_{\mathrm{FO}}}
\newcommand{\RCTF}{R_{\mathrm{CTF}}}
\newcommand{\Rzero}{R_{0}}

\title{Proofs of the Fourier-optics envelope identities\\
in aberration-corrected LEEM}
\author{Wilson Lau}
\date{24 August 2026}

\begin{document}
\maketitle

\begin{abstract}
We give complete analytic proofs of the appendix identities of
Yu, Lau and Altman,
\emph{Ultramicroscopy} \textbf{200} (2019) 160--168,
for Fourier-optics (FO) image formation in aberration-corrected
low-energy electron microscopy (ac-LEEM).
After recording the modelling approximations that factor the bilinear
kernel, we prove: the Gaussian characteristic functions that produce the
spatial envelope~\eqref{eq:paper5};
the complex Gaussian integral that produces the chromatic
envelope~\eqref{eq:paper6a}--\eqref{eq:paper6b};
the polar form of the chromatic prefactor;
the Gaussian algebra of the FO/CTF ratios~\eqref{eq:paper7};
and the Jacobi--Anger expansion of a sinusoidal phase object.
All statements are identities under the hypotheses named in the text.
The first-order tilt truncation, the dropping of mixed tilt--energy terms,
and the frozen-aperture approximation are \emph{not} identities and are
labelled as such.
These proofs match the Lean~4 formalization in this repository.
\end{abstract}

\tableofcontents
\newpage

\section{Introduction}
\label{sec:intro}

Image formation in LEEM is bilinear in the object spectrum:
partially coherent illumination mixes spatial frequencies
$(\mathbf{q},\mathbf{q}')$ rather than multiplying a single contrast
transfer function (CTF).
Yu, Lau and Altman~\cite{Yu2019} write the ensemble-averaged kernel as an
integral of a coherent kernel against a Gaussian source density and a
Gaussian energy density, then evaluate those integrals in Appendix~A1.

This note proves every closed form used in that appendix and in the
subsequent FO/CTF comparison, with the spatial-frequency convention of
Tromp--Schramm~\cite{Tromp2012,Schramm2012}:
the wave aberration $\chi$ is measured in \emph{waves}, the transfer
factor is $W=\ee^{2\pi\ii\chi}$, and $\lambda$ is restored so that $q$ is a
true spatial frequency.
Object phase is in radians and is independent of that $2\pi$ convention.
These identities are the forward kernel required to invert a
through-focal intensity series for the exit wave; the inverse against
this kernel is treated in a companion note.

\paragraph{What is proved.}
Gaussian Fourier transforms; FWHM--variance conversion; the gradient of
$\chi_S$; the exact splitting $\Delta\chi_C=b_1\varepsilon+b_2\varepsilon^2$;
closed forms of $E_S$ and $\ECtot$; polar form of $\ECC$;
$\Gamma_S=\ee^{-2\kappa\,u v}$; axis and perfect-coherence identities;
$\RCTF=\RFO\,\Gamma_C\Gamma_S$; Jacobi--Anger.

\paragraph{What is not proved.}
Experimental claims, numerical accuracy of Table~1 parameters, and the
modelling truncations of~\cite[Eqs.~(3)--(4)]{Yu2019}.

\section{Notation and standing conventions}
\label{sec:notation}

Vectors are bold; $q=\abs{\mathbf{q}}$.
In one dimension, $q,q',k$ are \emph{signed} scalars, so $qq'$ is a product
of signed reals.
The Fourier kernel is the unitary $2\pi$-in-the-exponent convention
\[
  \ee^{2\pi\ii\,\mathbf{q}\cdot\mathbf{r}},\qquad
  \dd^2q=\dd q_x\,\dd q_y.
\]
The principal branch of the complex logarithm has argument in
$(-\pi,\pi]$, and $z^{1/2}=\exp(\tfrac12\mathrm{Log}\,z)$ with
$z^{1/2}>0$ for $z>0$ real.
The same branch defines $z^{-1/2}$.

\begin{definition}[Object wave]
\label{def:object}
$\psi_0(\mathbf{r})=\sigma(\mathbf{r})\,\ee^{\ii\phi(\mathbf{r})}$,
with $\phi$ in radians.
\end{definition}

\begin{definition}[Spatial frequency]
\label{def:q}
$\mathbf{q}=\boldsymbol{\alpha}/\lambda$, where $\boldsymbol{\alpha}$ is the emission-angle
vector from the virtual object and $\lambda$ is the wavelength after
acceleration to column energy $E$.
\end{definition}

\begin{definition}[Aperture]
\label{def:aperture}
For a circular contrast aperture of cutoff $\qap=\alpha_{\mathrm{ap}}/\lambda$,
\[
  M(\mathbf{q})=\mathbf{1}_{\{\abs{\mathbf{q}}\le\qap\}}.
\]
In particular $M$ is real, so $M^*(\mathbf{q}')=M(\mathbf{q}')$.
\end{definition}

\begin{approximation}[Frozen aperture]
\label{approx:aperture}
For a small source, $M(\mathbf{q}+\mathbf{k})\approx M(\mathbf{q})$.
This is a modelling hypothesis, not an identity.
\end{approximation}

\begin{definition}[Spherical and defocus aberration]
\label{def:chiS}
In waves,
\begin{equation}
\label{eq:chiS}
\chi_S(\mathbf{q},\Dz)
=\tfrac14 C_3\lambda^3 q^4
+\tfrac16 C_5\lambda^5 q^6
+\tfrac12\Dz\,\lambda\,q^2,
\end{equation}
and $W_S=\ee^{2\pi\ii\chi_S}$.
The sign of $\Dz$ is the Tromp--Schramm convention: $\Dz>0$ adds
positive optical path for off-axis rays.
Non-corrected (nac) LEEM keeps $C_3$ and drops $C_5$;
aberration-corrected (ac) LEEM sets $C_3=0$ and keeps $C_5$.
In one dimension replace $q^{2n}$ by the even power of the signed scalar.
\end{definition}

\begin{definition}[Chromatic aberration]
\label{def:chiC}
\begin{equation}
\label{eq:chiC}
\chi_C(\mathbf{q},\varepsilon)
=\tfrac12 C_C\lambda\Bigl(\frac{\varepsilon}{E}\Bigr)q^2
+\tfrac12 C_{CC}\lambda\Bigl(\frac{\varepsilon}{E}\Bigr)^2 q^2
+\tfrac14 C_{3C}\lambda^3\Bigl(\frac{\varepsilon}{E}\Bigr)q^4,
\end{equation}
and $W_C=\ee^{2\pi\ii\chi_C}$.
Nac keeps $C_C$ and drops $C_{CC},C_{3C}$;
ac sets $C_C=0$ and keeps $C_{CC},C_{3C}$.
\end{definition}

\begin{definition}[Coherent kernel]
\label{def:R0}
Paper Eq.~(2):
\begin{equation}
\label{eq:R0}
\Rzero(\mathbf{q},\mathbf{q}',\Dz,\varepsilon)
=M(\mathbf{q})M(\mathbf{q}')
\,W_S(\mathbf{q},\Dz)\,W_S^*(\mathbf{q}',\Dz)
\,W_C(\mathbf{q},\varepsilon)\,W_C^*(\mathbf{q}',\varepsilon).
\end{equation}
\end{definition}

\begin{lemma}[Unitary transfer]
\label{lem:WW}
If $\chi$ is real then
$W(\mathbf{q})W^*(\mathbf{q}')=\ee^{2\pi\ii(\chi(\mathbf{q})-\chi(\mathbf{q}'))}$.
At $\varepsilon=0$, $\chi_C(\cdot,0)=0$, so $W_C(\cdot,0)=1$.
\end{lemma}
\begin{proof}
$W=\ee^{2\pi\ii\chi}$ lies on the unit circle, hence $W^{-1}=W^*$.
The second claim is $\chi_C(\mathbf{q},0)=0$ by~\eqref{eq:chiC}.
\end{proof}

\begin{definition}[Ensemble kernel]
\label{def:R}
Paper Eq.~(1):
\begin{equation}
\label{eq:R}
R(\mathbf{q},\mathbf{q}',\Dz)
=\int_{\R^2}\dd^2k\int_{\R}\dd\varepsilon\;
s(\mathbf{k})\,c(\varepsilon)\,
\Rzero(\mathbf{q}+\mathbf{k},\mathbf{q}'+\mathbf{k},\Dz,\varepsilon).
\end{equation}
\end{definition}

\begin{approximation}[First-order tilt, no mixed $k\varepsilon$]
\label{approx:taylor}
Expand $\chi_S(\mathbf{q}+\mathbf{k})$ to first order in $\mathbf{k}$ and
drop mixed tilt--energy terms in $\chi_C(\mathbf{q}+\mathbf{k},\varepsilon)$.
Together with~\cref{approx:aperture} this is paper Eq.~(4):
\begin{equation}
\label{eq:R0approx}
\begin{aligned}
\Rzero(\mathbf{q}+\mathbf{k},\mathbf{q}'+\mathbf{k},\Dz,\varepsilon)
&\approx
\Rzero(\mathbf{q},\mathbf{q}',\Dz,0)
\,\ee^{2\pi\ii\,\mathbf{a}\cdot\mathbf{k}}
\,\ee^{2\pi\ii(b_1\varepsilon+b_2\varepsilon^2)},
\end{aligned}
\end{equation}
with $\mathbf{a}$, $b_1$, $b_2$ as in~\cref{sec:aberration}.
The $(\mathbf{k},\varepsilon)$ integral then factors as
$\EStot\,\ECtot$.
The remainder of this note evaluates those two integrals exactly.
\end{approximation}

\section{Gaussian calculus}
\label{sec:gaussian}

\begin{definition}[Normalised Gaussians]
\label{def:gaussians}
For any width $\sigma>0$ write
\[
g_\sigma(x)
=\frac{1}{\sigma\sqrt{2\pi}}
\exp\Bigl(-\frac{x^2}{2\sigma^2}\Bigr).
\]
The energy and source densities of the paper are
\begin{align}
c(\varepsilon)
&=g_{\sE}(\varepsilon),
\label{eq:c}\\
s_{1\mathrm{D}}(k)
&=g_{\sill}(k),
\qquad
s(\mathbf{k})
=g_{\sill}(k_x)\,g_{\sill}(k_y)
=\frac{1}{2\pi\sill^2}
\exp\Bigl(-\frac{\abs{\mathbf{k}}^2}{2\sill^2}\Bigr).
\label{eq:s}
\end{align}
\end{definition}

\begin{lemma}[FWHM of a centred Gaussian]
\label{lem:fwhm}
Let $f(x)=\exp(-x^2/(2\sigma^2))$ with $\sigma>0$.
If $f(\pm\Delta/2)=\tfrac12 f(0)$, then
\begin{equation}
\label{eq:fwhm}
\Delta=2\sqrt{2\ln 2}\,\sigma,
\qquad
\sigma^2=\frac{\Delta^2}{8\ln 2}.
\end{equation}
\end{lemma}
\begin{proof}
$f(0)=1$, so the half-maximum condition is
$\exp\bigl(-(\Delta/2)^2/(2\sigma^2)\bigr)=1/2$.
Taking $-\log$ yields $(\Delta/2)^2/(2\sigma^2)=\ln 2$, hence
$\Delta^2/(8\sigma^2)=\ln 2$.
The positive root is~\eqref{eq:fwhm}.
Squaring $\Delta=2\sqrt{2\ln 2}\,\sigma$ recovers the variance formula.
\end{proof}

The paper identifies $\DE$ with the FWHM of $c$ and $\qill=\alpha_{\mathrm{ill}}/\lambda$
with the isotropic FWHM of $s$.
Write $F:=2\sqrt{2\ln 2}$ for the FWHM factor.

\begin{lemma}[Product structure]
\label{lem:product}
$s(\mathbf{k})=s_{1\mathrm{D}}(k_x)\,s_{1\mathrm{D}}(k_y)$.
\end{lemma}
\begin{proof}
$\abs{\mathbf{k}}^2=k_x^2+k_y^2$ splits the exponential.
The prefactors multiply:
$(1/(\sill\sqrt{2\pi}))^2=1/(2\pi\sill^2)$ because
$\sqrt{2\pi}\,\sqrt{2\pi}=2\pi$.
\end{proof}

\begin{lemma}[Complex Gaussian integral]
\label{lem:cgaussian}
Let $B,C,D\in\C$ with $\ReOp(B)<0$. Then
\begin{equation}
\label{eq:cgaussian}
\int_{\R}\exp(Bx^2+Cx+D)\,\dd x
=\Bigl(\frac{\pi}{-B}\Bigr)^{1/2}
\exp\Bigl(D-\frac{C^2}{4B}\Bigr),
\end{equation}
using the principal square root.
\end{lemma}
\begin{proof}
Complete the square:
$Bx^2+Cx+D=B\bigl(x+C/(2B)\bigr)^2+D-C^2/(4B)$.
The hypothesis $\ReOp(B)<0$ gives Gaussian decay in a neighbourhood of the
real axis, so the contour may be shifted and the integral equals
$\exp(D-C^2/(4B))$ times $\int\exp(Bw^2)\,\dd w$.
The standard evaluation of the latter is $(\pi/(-B))^{1/2}$.
Because $\ReOp(-B)>0$, the number $-B$ lies in the open right half-plane,
so $\Arg(-B)\in(-\pi/2,\pi/2)$ and the principal square root never meets
its branch cut (the negative reals).
When $B$ is real and negative this recovers the positive root
$\sqrt{\pi/\abs{B}}$.
\end{proof}

\begin{theorem}[One-dimensional characteristic function]
\label{thm:char1d}
For $\sigma>0$ and $a\in\R$,
\begin{equation}
\label{eq:char1d}
\int_{\R}g_\sigma(k)\,\ee^{2\pi\ii a k}\,\dd k
=\exp(-2\pi^2\sigma^2 a^2).
\end{equation}
In particular this holds for $s_{1\mathrm{D}}=g_{\sill}$ and for $c=g_{\sE}$.
\end{theorem}
\begin{proof}
The integrand is a constant times $\exp(Bk^2+Ck)$ with
$B=-1/(2\sigma^2)$ (real) and $C=2\pi\ii a$.
Then $\ReOp(B)<0$.
\Cref{lem:cgaussian} gives the prefactor
\[
\frac{1}{\sigma\sqrt{2\pi}}
\Bigl(\frac{\pi}{1/(2\sigma^2)}\Bigr)^{1/2}
=\frac{1}{\sigma\sqrt{2\pi}}\sqrt{2\pi\sigma^2}=1,
\]
using the positive square root ($\sigma>0$).
The exponent is $-C^2/(4B)$.
Now $C^2=(2\pi\ii a)^2=-4\pi^2 a^2$, so
\[
-\frac{C^2}{4B}=\frac{4\pi^2 a^2}{4B}=\frac{\pi^2 a^2}{B}
=\pi^2 a^2\cdot\bigl(-2\sigma^2\bigr)=-2\pi^2\sigma^2 a^2.
\]
Equivalently, the characteristic function of $\mathcal{N}(0,\sigma^2)$ is
$\mathbb{E}[\ee^{\ii t k}]=\ee^{-\sigma^2 t^2/2}$ at $t=2\pi a$.
\end{proof}

\begin{theorem}[Two-dimensional characteristic function]
\label{thm:char2d}
For $\sigma>0$ and $\mathbf{a}\in\R^2$, write
$s_\sigma(\mathbf{k})=g_\sigma(k_x)g_\sigma(k_y)$. Then
\begin{equation}
\label{eq:char2d}
\int_{\R^2}s_\sigma(\mathbf{k})\,\ee^{2\pi\ii\,\mathbf{a}\cdot\mathbf{k}}\,\dd^2k
=\exp(-2\pi^2\sigma^2\abs{\mathbf{a}}^2).
\end{equation}
The illumination density $s$ is the case $\sigma=\sill$.
\end{theorem}
\begin{proof}
The density factors as a product of one-dimensional Gaussians of the same
width, and
$\mathbf{a}\cdot\mathbf{k}=a_x k_x+a_y k_y$ factors the exponential.
The double integral is a product of two copies of~\cref{thm:char1d},
and the exponents add.
\end{proof}

\begin{corollary}[FWHM packaging]
\label{cor:kappa}
If $\qill=F\sill$ and
$\kappa:=\pi^2\qill^2/(4\ln 2)$, then $2\pi^2\sill^2=\kappa$.
Likewise $2\pi^2\sE^2=\pi^2\DE^2/(4\ln 2)$.
\end{corollary}
\begin{proof}
Substitute $\sigma^2=\Delta^2/(8\ln 2)$ from~\cref{lem:fwhm} into
$2\pi^2\sigma^2$:
\[
2\pi^2\cdot\frac{\Delta^2}{8\ln 2}=\frac{\pi^2\Delta^2}{4\ln 2}.
\]
\end{proof}

\section{Aberration identities}
\label{sec:aberration}

\begin{theorem}[Gradient of $\chi_S$]
\label{thm:grad}
Viewing~\eqref{eq:chiS} as a radial $C^\infty$ function on $\R^2$,
\begin{equation}
\label{eq:grad}
\nabla_{\mathbf{q}}\chi_S(\mathbf{q},\Dz)
=\bigl(C_3\lambda^3 q^2+C_5\lambda^5 q^4+\Dz\,\lambda\bigr)\,\mathbf{q}.
\end{equation}
In one dimension,
$\partial_q\chi_S=C_3\lambda^3 q^3+C_5\lambda^5 q^5+\Dz\,\lambda\,q$.
\end{theorem}
\begin{proof}
Write $\rho=q^2=\abs{\mathbf{q}}^2$, so
$\chi_S=\tfrac14 C_3\lambda^3\rho^2+\tfrac16 C_5\lambda^5\rho^3+\tfrac12\Dz\lambda\,\rho$.
The Euclidean gradient is $\nabla\rho=2\mathbf{q}$.
Chain rule:
$\nabla(\rho^2)=4\rho\,\mathbf{q}=4q^2\mathbf{q}$ and
$\nabla(\rho^3)=6\rho^2\mathbf{q}=6q^4\mathbf{q}$.
The numerical coefficients cancel: $\tfrac14\cdot 4=1$ and $\tfrac16\cdot 6=1$
and $\tfrac12\cdot 2=1$.
In one dimension, differentiate the even powers $q^4$, $q^6$, $q^2$
directly; odd powers keep the sign of $q$.
\end{proof}

\begin{definition}[Tilt vector]
\label{def:a}
\begin{equation}
\label{eq:a}
\mathbf{a}(\mathbf{q},\mathbf{q}',\Dz)
:=\nabla\chi_S(\mathbf{q},\Dz)-\nabla\chi_S(\mathbf{q}',\Dz).
\end{equation}
In one dimension,
$a=C_3\lambda^3(q^3-{q'}^3)+C_5\lambda^5(q^5-{q'}^5)+\Dz\lambda(q-q')$.
On the diagonal, $\mathbf{a}(\mathbf{q},\mathbf{q})=0$.
On the axis $\mathbf{q}'=\mathbf{0}$, $\nabla\chi_S(\mathbf{0})=0$, so
$\mathbf{a}(\mathbf{q},\mathbf{0})=\nabla\chi_S(\mathbf{q})$.
\end{definition}

\begin{theorem}[First-order remainder]
\label{thm:taylor}
For each fixed $\mathbf{q},\Dz$,
\[
\chi_S(\mathbf{q}+\mathbf{k},\Dz)-\chi_S(\mathbf{q},\Dz)
-\mathbf{k}\cdot\nabla\chi_S(\mathbf{q},\Dz)
=o(\abs{\mathbf{k}})\qquad(\mathbf{k}\to\mathbf{0}).
\]
\end{theorem}
\begin{proof}
This is the definition of differentiability.
Since $\chi_S$ is a polynomial the remainder is in fact $O(\abs{\mathbf{k}}^2)$.
The little-o statement is what paper Eq.~(3) uses; it does \emph{not}
justify dropping the remainder uniformly in $\mathbf{q}$, which remains
\cref{approx:taylor}.
\end{proof}

\begin{theorem}[$b_1$/$b_2$ splitting]
\label{thm:b1b2}
If $E\neq 0$, then
$\chi_C(\mathbf{q},\varepsilon)-\chi_C(\mathbf{q}',\varepsilon)
=b_1\varepsilon+b_2\varepsilon^2$ with
\begin{equation}
\label{eq:b1b2}
\begin{aligned}
b_1(\mathbf{q},\mathbf{q}')
&=\frac{C_C\lambda}{2E}(q^2-{q'}^2)
+\frac{C_{3C}\lambda^3}{4E}(q^4-{q'}^4),\\
b_2(\mathbf{q},\mathbf{q}')
&=\frac{C_{CC}\lambda}{2E^2}(q^2-{q'}^2).
\end{aligned}
\end{equation}
In two dimensions, $q^2$ means $\abs{\mathbf{q}}^2$.
Units: $b_1$ is waves/energy, $b_2$ is waves/energy$^2$.
\end{theorem}
\begin{proof}
Subtract~\eqref{eq:chiC} at $\mathbf{q}$ and $\mathbf{q}'$ and collect
powers of $\varepsilon$.
The $C_C$ term contributes $\tfrac12 C_C\lambda E^{-1}(q^2-{q'}^2)\varepsilon$;
the $C_{3C}$ term contributes $\tfrac14 C_{3C}\lambda^3 E^{-1}(q^4-{q'}^4)\varepsilon$;
the $C_{CC}$ term contributes $\tfrac12 C_{CC}\lambda E^{-2}(q^2-{q'}^2)\varepsilon^2$.
No other powers appear.
This is an identity, not an approximation.
\end{proof}

\begin{corollary}[Specialisations]
\label{cor:b1b2}
On the diagonal, $b_1(\mathbf{q},\mathbf{q})=b_2(\mathbf{q},\mathbf{q})=0$.
Antisymmetry: $b_j(\mathbf{q},\mathbf{q}')=-b_j(\mathbf{q}',\mathbf{q})$.
Under ac ($C_C=0$),
$b_1=(C_{3C}\lambda^3/(4E))(q^4-{q'}^4)$.
Under nac ($C_{CC}=C_{3C}=0$), $b_2=0$ and
$b_1=(C_C\lambda/(2E))(q^2-{q'}^2)$.
\end{corollary}
\begin{proof}
Immediate from~\eqref{eq:b1b2}.
\end{proof}

\section{Spatial envelope}
\label{sec:spatial}

\begin{definition}[Spatial envelope integral]
\label{def:ES}
After~\cref{approx:taylor},
\begin{equation}
\label{eq:ESint}
\EStot(\mathbf{q},\mathbf{q}',\Dz)
:=\int_{\R^2}
s(\mathbf{k})\,\ee^{2\pi\ii\,\mathbf{a}\cdot\mathbf{k}}\,\dd^2k,
\end{equation}
and likewise in one dimension with $s_{1\mathrm{D}}$ and scalar $a$.
\end{definition}

\begin{theorem}[Closed form in $\sill$]
\label{thm:ES-sigma}
If $\sill>0$, then
$\EStot=\exp(-2\pi^2\sill^2\abs{\mathbf{a}}^2)$.
The envelope is real, even under
$(\mathbf{q},\mathbf{q}')\leftrightarrow(\mathbf{q}',\mathbf{q})$,
and equals $1$ on the diagonal.
\end{theorem}
\begin{proof}
Apply~\cref{thm:char2d} (or~\cref{thm:char1d} in one dimension).
Reality: the exponent is real.
Symmetry: $\mathbf{a}(\mathbf{q}',\mathbf{q})=-\mathbf{a}(\mathbf{q},\mathbf{q}')$,
and $\abs{\mathbf{a}}^2$ is even.
On the diagonal $\mathbf{a}=\mathbf{0}$.
\end{proof}

\begin{theorem}[Paper Eq.~(5)]
\label{thm:ES-fwhm}
If $\qill>0$ and $\kappa=\pi^2\qill^2/(4\ln 2)$, then
\begin{equation}
\label{eq:paper5}
\boxed{
\EStot(\mathbf{q},\mathbf{q}',\Dz)
=\exp\Bigl(-\frac{\pi^2\qill^2}{4\ln 2}\,
\abs{\mathbf{a}(\mathbf{q},\mathbf{q}',\Dz)}^2\Bigr)
=\ee^{-\kappa\abs{\mathbf{a}}^2}.
}
\end{equation}
\end{theorem}
\begin{proof}
Combine~\cref{thm:ES-sigma} with~\cref{cor:kappa}.
\end{proof}

\begin{corollary}[Expanded $\abs{\mathbf{a}}^2$]
\label{cor:a2}
Let $u(\mathbf{q}):=C_3\lambda^3 q^2+C_5\lambda^5 q^4+\Dz\lambda$,
so $\nabla\chi_S=u(\mathbf{q})\,\mathbf{q}$ and
$\mathbf{a}=u(\mathbf{q})\mathbf{q}-u(\mathbf{q}')\mathbf{q}'$.
Then
\[
\abs{\mathbf{a}}^2
=u(\mathbf{q})^2 q^2+u(\mathbf{q}')^2{q'}^2
-2\,u(\mathbf{q})u(\mathbf{q}')\,(\mathbf{q}\cdot\mathbf{q}').
\]
The inner product $\mathbf{q}\cdot\mathbf{q}'$ cannot be replaced by
$q\,q'$.
\end{corollary}
\begin{proof}
Expand $\abs{\mathbf{x}-\mathbf{y}}^2$.
\end{proof}

\begin{corollary}[One-dimensional polynomial form]
\label{cor:ES1d}
In one dimension~\eqref{eq:paper5} is
\[
\EStot(q,q',\Dz)
=\exp\Biggl(
-\kappa
\Bigl[
C_3\lambda^3(q^3-{q'}^3)
+C_5\lambda^5(q^5-{q'}^5)
+\Dz\lambda(q-q')
\Bigr]^2
\Biggr).
\]
One may factor $(q-q')^2$ via $q^3-{q'}^3=(q-q')(q^2+qq'+{q'}^2)$.
For ac, $C_3=0$.
\end{corollary}

\begin{corollary}[CTF slice]
\label{cor:ES-ctf}
$\EStot(\mathbf{q},\mathbf{0},\Dz)=\exp(-\kappa\abs{\nabla\chi_S(\mathbf{q})}^2)$,
the separable CTF spatial envelope.
\end{corollary}
\begin{proof}
\Cref{def:a}: $\mathbf{a}(\mathbf{q},\mathbf{0})=\nabla\chi_S(\mathbf{q})$.
\end{proof}

\section{Chromatic envelope}
\label{sec:chromatic}

\begin{definition}[Chromatic envelope integral]
\label{def:EC}
\begin{equation}
\label{eq:ECint}
\ECtot(\mathbf{q},\mathbf{q}')
:=\int_{\R}
c(\varepsilon)\,
\exp\bigl(2\pi\ii(b_1\varepsilon+b_2\varepsilon^2)\bigr)
\,\dd\varepsilon.
\end{equation}
\end{definition}

\begin{lemma}[Quadratic coefficient]
\label{lem:B}
The integrand of~\eqref{eq:ECint} is a constant times
$\exp(B\varepsilon^2+C\varepsilon)$ with
\begin{equation}
\label{eq:B}
B=-\frac{1}{2\sE^2}+2\pi\ii b_2,
\qquad
C=2\pi\ii b_1.
\end{equation}
If $\sE>0$ then $\ReOp(B)=-1/(2\sE^2)<0$, independently of $b_2$.
Writing $z:=1-4\pi\ii b_2\sE^2$ (so $\ReOp(z)=1>0$), one has
$B=-z/(2\sE^2)$.
\end{lemma}
\begin{proof}
Add the Gaussian exponent $-\varepsilon^2/(2\sE^2)$ to the phase
$2\pi\ii(b_1\varepsilon+b_2\varepsilon^2)$.
The imaginary part of $B$ is $2\pi b_2$ and does not affect $\ReOp(B)$.
The identity $B=-z/(2\sE^2)$ is algebra:
\[
-\frac{z}{2\sE^2}
=-\frac{1-4\pi\ii b_2\sE^2}{2\sE^2}
=-\frac{1}{2\sE^2}+2\pi\ii b_2.
\]
Thus~\cref{lem:cgaussian} applies for every real $b_2$ as soon as
$\sE>0$.
\end{proof}

\begin{lemma}[Square-root identity]
\label{lem:sqrt-div}
If $A>0$ is real and $D\in\C\setminus\{0\}$ has $\Arg(D)\neq\pi$, then
$(A/D)^{1/2}=A^{1/2}D^{-1/2}$ on principal branches.
This applies to $D=z=1-4\pi\ii b_2\sE^2$, because $\ReOp(z)=1$ implies
$\Arg(z)\in(-\pi/2,\pi/2)$.
\end{lemma}
\begin{proof}
$\mathrm{Log}(A/D)=\log A-\mathrm{Log}\,D$ provided $\Arg(D)\neq\pi$
(the identity $\mathrm{Log}(1/w)=-\mathrm{Log}\,w$ fails on the negative
reals).
Exponentiate $\tfrac12\mathrm{Log}(A/D)$.
\end{proof}

\begin{theorem}[Paper Eq.~(6b)]
\label{thm:ECC}
Define
\begin{equation}
\label{eq:paper6b}
\boxed{
\ECC(\mathbf{q},\mathbf{q}')
:=\bigl(1-4\pi\ii\,b_2\sE^2\bigr)^{-1/2}
=z^{-1/2}
}
\end{equation}
on the principal branch, with value $1$ at $b_2=0$.
Then the Gaussian-integral prefactor of~\eqref{eq:ECint} equals $\ECC$,
and $\ECC^2=z^{-1}$.
\end{theorem}
\begin{proof}
\Cref{lem:cgaussian} produces the prefactor
\[
\frac{1}{\sE\sqrt{2\pi}}
\Bigl(\frac{\pi}{-B}\Bigr)^{1/2}.
\]
From~\cref{lem:B}, $\pi/(-B)=\pi\cdot(2\sE^2)/z=2\pi\sE^2/z$.
\Cref{lem:sqrt-div} with $A=2\pi\sE^2>0$ gives
$(\pi/(-B))^{1/2}=(2\pi\sE^2)^{1/2}z^{-1/2}$.
The positive real root $(2\pi\sE^2)^{1/2}=\sE\sqrt{2\pi}$ cancels the
Gaussian normalisation, leaving $z^{-1/2}=\ECC$.
At $b_2=0$, $z=1$ and $1^{-1/2}=1$ (not $-1$).
The identity $\ECC^2=z^{-1}$ is
$z^{-1/2}z^{-1/2}=z^{-1}$, valid because $\Arg(z)\neq\pi$.
\end{proof}

In FWHM variables,
$4\pi b_2\sE^2=\pi b_2\DE^2/(2\ln 2)$, so
$z=1-\ii\,(\pi b_2\DE^2)/(2\ln 2)$.

\begin{theorem}[Paper Eq.~(6a)]
\label{thm:EC-closed}
If $\sE>0$, then
\begin{equation}
\label{eq:paper6a}
\boxed{
\ECtot
=\ECC\exp\bigl(-2\pi^2\sE^2 b_1^2\,\ECC^2\bigr)
=\bigl(1-4\pi\ii b_2\sE^2\bigr)^{-1/2}
\exp\Bigl(
-\frac{2\pi^2\sE^2 b_1^2}{1-4\pi\ii b_2\sE^2}
\Bigr).
}
\end{equation}
\end{theorem}
\begin{proof}
\Cref{lem:cgaussian} supplies the exponent $-C^2/(4B)$ with
$C=2\pi\ii b_1$.
Then $C^2=(2\pi\ii b_1)^2=-4\pi^2 b_1^2$, hence
\[
-\frac{C^2}{4B}=\frac{4\pi^2 b_1^2}{4B}=\frac{\pi^2 b_1^2}{B}.
\]
Substitute $B=-z/(2\sE^2)$:
\[
\frac{\pi^2 b_1^2}{B}
=\pi^2 b_1^2\cdot\frac{-2\sE^2}{z}
=-2\pi^2\sE^2 b_1^2/z.
\]
Since $\ECC^2=z^{-1}$, this is $-2\pi^2\sE^2 b_1^2\,\ECC^2$.
Multiply by the prefactor of~\cref{thm:ECC}.
\end{proof}

\begin{remark}[Hanßen--Trepte packaging]
\label{rem:16ln2}
One has
$2\pi^2\sE^2 b_1^2=(\DE)^2(2\pi b_1)^2/(16\ln 2)$,
where $2\pi b_1$ is the coefficient of $\varepsilon$ in the
\emph{radian} phase $2\pi\Delta\chi_C$.
Both $\pi$ (from $W=\ee^{2\pi\ii\chi}$) and $16\ln 2$ (from FWHM) must appear.
\end{remark}

\begin{corollary}[Nac / $b_2=0$]
\label{cor:nac}
If $b_2=0$ then $\ECC=1$ and
$\ECtot=\exp(-2\pi^2\sE^2 b_1^2)$.
On the nac CTF slice $\mathbf{q}'=\mathbf{0}$,
$b_1=C_C\lambda q^2/(2E)$, hence
\begin{equation}
\label{eq:nac-ctf}
\ECtot(\mathbf{q},\mathbf{0})
=\exp\Biggl(
-\frac{\bigl(\pi C_C\lambda\DE\,q^2\bigr)^2}{16\ln 2\,E^2}
\Biggr),
\end{equation}
the standard chromatic envelope of Schramm \emph{et al.}~\cite{Schramm2012}.
\end{corollary}
\begin{proof}
\Cref{thm:EC-closed} with $z=1$, then substitute $b_1$ and~\cref{cor:kappa}:
\[
2\pi^2\sE^2 b_1^2
=\frac{\pi^2\DE^2}{4\ln 2}\cdot\frac{C_C^2\lambda^2 q^4}{4E^2}
=\frac{(\pi C_C\lambda\DE q^2)^2}{16\ln 2\,E^2}.
\]
\end{proof}

\begin{corollary}[Conjugation]
\label{cor:ECconj}
$\ECtot(-b_1,-b_2)=\conj{\ECtot(b_1,b_2)}$.
Equivalently $\ECtot(\mathbf{q}',\mathbf{q})=\conj{\ECtot(\mathbf{q},\mathbf{q}')}$.
\end{corollary}
\begin{proof}
The density $c$ is real, and
$\exp(2\pi\ii\psi)^*=\exp(-2\pi\ii\psi)$ for real $\psi$.
Thus the integrand at $(-b_1,-b_2)$ is the conjugate of the integrand at
$(b_1,b_2)$.
Integrate, or apply~\cref{cor:b1b2} to swap $(\mathbf{q},\mathbf{q}')$.
\end{proof}

\section{Polar form of the chromatic envelope}
\label{sec:polar}

Write $y:=4\pi b_2\sE^2=\pi b_2\DE^2/(2\ln 2)$, so
$\ECC=(1-\ii y)^{-1/2}$.
Under ac,
$y=\gamma(q^2-{q'}^2)$ with
\begin{equation}
\label{eq:gamma}
\gamma:=\frac{\pi C_{CC}\lambda\DE^2}{4\ln 2\,E^2}.
\end{equation}

\begin{lemma}[Modulus and argument of $1-\ii y$]
\label{lem:arg}
For $y\in\R$, $\abs{1-\ii y}=\sqrt{1+y^2}$ and
$\Arg(1-\ii y)=-\arctan y$.
In particular $\Arg(1-\ii y)\in(-\pi/2,\pi/2)$.
\end{lemma}
\begin{proof}
Real part $1>0$, imaginary part $-y$.
The modulus is $\sqrt{1+y^2}$.
For a number in the closed right half-plane,
$\Arg=\mathrm{atan2}(\ImOp,\ReOp)=\mathrm{atan2}(-y,1)=-\arctan y$,
where $\arctan:\R\to(-\pi/2,\pi/2)$.
Because the real part is identically $1$, the argument never jumps.
\end{proof}

\begin{theorem}[Polar form of $\ECC$]
\label{thm:ECC-polar}
On the principal branch,
\begin{equation}
\label{eq:ECC-polar}
(1-\ii y)^{-1/2}
=(1+y^2)^{-1/4}
\exp\bigl(\tfrac{\ii}{2}\arctan y\bigr).
\end{equation}
Hence under ac,
\[
\ECC
=\bigl[1+\gamma^2(q^2-{q'}^2)^2\bigr]^{-1/4}
\exp\Bigl(\tfrac{\ii}{2}\arctan\bigl(\gamma(q^2-{q'}^2)\bigr)\Bigr).
\]
\end{theorem}
\begin{proof}
$z^{-1/2}=\abs{z}^{-1/2}\exp(-\ii\Arg(z)/2)$.
By~\cref{lem:arg}, $\abs{z}^{-1/2}=(1+y^2)^{-1/4}$ and
$-\Arg(z)/2=\tfrac12\arctan y$.
The phase of $\ECC$ therefore lies in $(-\pi/4,\pi/4)$ and never jumps.
At $y=0$ one recovers $+1$.
\end{proof}

\begin{lemma}[Polar form of $1/z$]
\label{lem:1z}
$1/(1-\ii y)=(1+\ii y)/(1+y^2)$.
\end{lemma}
\begin{proof}
Multiply numerator and denominator by $1+\ii y$.
The denominator is $\abs{1-\ii y}^2=1+y^2$.
\end{proof}

\begin{theorem}[Polar form of the remaining exponential, ac]
\label{thm:exp-polar}
Let
\begin{equation}
\label{eq:eta}
\eta:=\frac{(\pi C_{3C}\lambda^3\DE)^2}{64\ln 2\,E^2},
\end{equation}
so that $2\pi^2\sE^2 b_1^2=\eta(q^4-{q'}^4)^2$ under ac.
Then
\begin{equation}
\label{eq:exp-polar}
\begin{aligned}
\exp(-2\pi^2\sE^2 b_1^2\,\ECC^2)
&=\exp\Biggl(
-\frac{\eta(q^4-{q'}^4)^2}{1+\gamma^2(q^2-{q'}^2)^2}
\Biggr)\\
&\qquad\cdot
\exp\Biggl(
-\ii\,
\frac{\eta\gamma(q^2-{q'}^2)(q^4-{q'}^4)^2}{1+\gamma^2(q^2-{q'}^2)^2}
\Biggr).
\end{aligned}
\end{equation}
\end{theorem}
\begin{proof}
Under ac, \cref{cor:b1b2} gives $b_1=(C_{3C}\lambda^3/(4E))(q^4-{q'}^4)$, so
\[
2\pi^2\sE^2 b_1^2
=\frac{\pi^2\DE^2}{4\ln 2}\cdot\frac{C_{3C}^2\lambda^6}{16E^2}(q^4-{q'}^4)^2
=\eta(q^4-{q'}^4)^2.
\]
Now $\ECC^2=1/z=1/(1-\ii y)$.
The exponent is
$-2\pi^2\sE^2 b_1^2\cdot(1+\ii y)/(1+y^2)$
by~\cref{lem:1z}.
The real part is $-\eta(q^4-{q'}^4)^2/(1+y^2)$;
the imaginary part is $-\eta(q^4-{q'}^4)^2 y/(1+y^2)$.
Substitute $y=\gamma(q^2-{q'}^2)$.
\end{proof}

\begin{theorem}[Amplitude and phase of $\ECtot$, ac]
\label{thm:A-theta}
Write $\ECtot=A\ee^{\ii\theta}$ with $A\ge 0$. Then
\begin{equation}
\label{eq:A-theta}
\begin{aligned}
A(q,q')
&=\bigl[1+\gamma^2(q^2-{q'}^2)^2\bigr]^{-1/4}
\exp\Biggl(
-\frac{\eta(q^4-{q'}^4)^2}{1+\gamma^2(q^2-{q'}^2)^2}
\Biggr),\\
\theta(q,q')
&=\tfrac12\arctan\bigl(\gamma(q^2-{q'}^2)\bigr)
-\frac{\eta\gamma(q^2-{q'}^2)(q^4-{q'}^4)^2}{1+\gamma^2(q^2-{q'}^2)^2}.
\end{aligned}
\end{equation}
Parity: $A(q,q')=A(q',q)=A(\abs{q},\abs{q'})$ and
$\theta(q,q')=-\theta(q',q)$.
\end{theorem}
\begin{proof}
Product of~\cref{thm:ECC-polar,thm:exp-polar}: moduli multiply, arguments add.
Parity follows because $y$ and $(q^4-{q'}^4)$ both change sign under
$(q,q')\leftrightarrow(q',q)$, while $y^2$ and $(q^4-{q'}^4)^2$ do not.
\end{proof}

\section{Envelope ratios and the CTF comparison}
\label{sec:ratios}

\begin{definition}[Paper Eqs.~(7a)--(7b)]
\label{def:gamma}
\begin{equation}
\label{eq:paper7}
\begin{aligned}
\Gamma_C(\mathbf{q},\mathbf{q}')
&=\frac{\ECtot(\mathbf{q},\mathbf{0})\,\ECtot(\mathbf{q}',\mathbf{0})}
{\ECtot(\mathbf{q},\mathbf{q}')},
\\
\Gamma_S(\mathbf{q},\mathbf{q}',\Dz)
&=\frac{\EStot(\mathbf{q},\mathbf{0},\Dz)\,\EStot(\mathbf{q}',\mathbf{0},\Dz)}
{\EStot(\mathbf{q},\mathbf{q}',\Dz)}.
\end{aligned}
\end{equation}
Printed~(7a) uses the \emph{non-conjugated} product
$\ECtot(\mathbf{q},\mathbf{0})\ECtot(\mathbf{q}',\mathbf{0})$.
For ac, $\ECtot$ is complex, so $\Gamma_C$ is complex; the plotted quantity
is $\abs{\Gamma_C}$.
\end{definition}

\begin{theorem}[$\Gamma_S$ algebra]
\label{thm:gammaS}
Let $\mathbf{u}=\nabla\chi_S(\mathbf{q},\Dz)$,
$\mathbf{v}=\nabla\chi_S(\mathbf{q}',\Dz)$, and
$\kappa=\pi^2\qill^2/(4\ln 2)$.
Then
\begin{equation}
\label{eq:gammaS}
\Gamma_S=\exp(-2\kappa\,\mathbf{u}\cdot\mathbf{v}).
\end{equation}
In one dimension,
$\Gamma_S=\exp\bigl(-2\kappa\,(\partial_q\chi_S(q))(\partial_q\chi_S(q'))\bigr)$.
\end{theorem}
\begin{proof}
$\EStot=\ee^{-\kappa\abs{\mathbf{a}}^2}$ with $\mathbf{a}=\mathbf{u}-\mathbf{v}$,
and $\nabla\chi_S(\mathbf{0})=\mathbf{0}$, so
$\EStot(\mathbf{q},\mathbf{0})=\ee^{-\kappa\abs{\mathbf{u}}^2}$ and
$\EStot(\mathbf{q}',\mathbf{0})=\ee^{-\kappa\abs{\mathbf{v}}^2}$.
Therefore
\[
\Gamma_S
=\exp\bigl(-\kappa\abs{\mathbf{u}}^2-\kappa\abs{\mathbf{v}}^2
+\kappa\abs{\mathbf{u}-\mathbf{v}}^2\bigr).
\]
Expand $\abs{\mathbf{u}-\mathbf{v}}^2=\abs{\mathbf{u}}^2+\abs{\mathbf{v}}^2-2\mathbf{u}\cdot\mathbf{v}$.
The quadratic terms cancel, leaving $\ee^{-2\kappa\,\mathbf{u}\cdot\mathbf{v}}$.
\end{proof}

\begin{corollary}[One-dimensional ac formula]
\label{cor:gammaS-ac}
If $C_3=0$ and $\qill>0$, then
\begin{equation}
\label{eq:gammaS-ac}
\Gamma_S(q,q',\Dz)
=\exp\Biggl[
-\frac{\pi^2\qill^2}{2\ln 2}
\bigl(C_5\lambda^5 q^5+\Dz\lambda\,q\bigr)
\bigl(C_5\lambda^5 {q'}^5+\Dz\lambda\,q'\bigr)
\Biggr].
\end{equation}
\end{corollary}
\begin{proof}
\Cref{thm:gammaS} with $2\kappa=\pi^2\qill^2/(2\ln 2)$ and
\cref{thm:grad} with $C_3=0$.
\end{proof}

\begin{theorem}[Perfect coherence]
\label{thm:perfect}
If $\qill=\DE=0$ (equivalently $\sill=\sE=0$), then every envelope equals $1$,
hence $\Gamma_C=\Gamma_S=1$ and $\Gamma_C\Gamma_S=1$ for all
$(\mathbf{q},\mathbf{q}',\Dz)$.
\end{theorem}
\begin{proof}
$\exp(0)=1$ and $1^{-1/2}=1$.
The closed forms extend by continuity to vanishing widths.
\end{proof}

\begin{theorem}[Axes]
\label{thm:axes}
$\Gamma_S(\mathbf{q},\mathbf{0},\Dz)=\Gamma_S(\mathbf{0},\mathbf{q}',\Dz)=1$.
$\Gamma_C(\mathbf{q},\mathbf{0})=1$.
If $\DE>0$ then $\abs{\Gamma_C(\mathbf{0},\mathbf{q}')}=1$.
\end{theorem}
\begin{proof}
Write $E_S(\mathbf{q},\mathbf{q}')$ for $\EStot(\mathbf{q},\mathbf{q}',\Dz)$.
Then $E_S(\mathbf{0},\mathbf{0})=1$ because $\mathbf{a}(\mathbf{0},\mathbf{0})=\mathbf{0}$.
Hence
\[
\Gamma_S(\mathbf{q},\mathbf{0},\Dz)
=\frac{E_S(\mathbf{q},\mathbf{0})\,E_S(\mathbf{0},\mathbf{0})}{E_S(\mathbf{q},\mathbf{0})}
=1,
\]
and likewise $\Gamma_S(\mathbf{0},\mathbf{q}',\Dz)=1$ by symmetry of $E_S$.

For the chromatic ratios, $b_1(\mathbf{0},\mathbf{0})=b_2(\mathbf{0},\mathbf{0})=0$,
so $\ECtot(\mathbf{0},\mathbf{0})=1$. Therefore
\[
\Gamma_C(\mathbf{q},\mathbf{0})
=\frac{\ECtot(\mathbf{q},\mathbf{0})\,\ECtot(\mathbf{0},\mathbf{0})}
{\ECtot(\mathbf{q},\mathbf{0})}
=1.
\]
For the other axis, $b_j(\mathbf{0},\mathbf{q}')=-b_j(\mathbf{q}',\mathbf{0})$, so
\cref{cor:ECconj} gives
$\ECtot(\mathbf{0},\mathbf{q}')=\conj{\ECtot(\mathbf{q}',\mathbf{0})}$.
Let $Z:=\ECtot(\mathbf{q}',\mathbf{0})$. Then $Z\neq 0$
($\ECC\neq 0$ and the exponential never vanishes), and
\[
\Gamma_C(\mathbf{0},\mathbf{q}')
=\frac{\ECtot(\mathbf{0},\mathbf{0})\,Z}{\conj{Z}}
=\frac{Z}{\conj{Z}},
\]
which has modulus $1$. (For ac this need not equal $1$: only the modulus is
identically one.)
\end{proof}

\begin{definition}[Working FO and CTF kernels]
\label{def:RFO}
After~\cref{approx:taylor},
\begin{align}
\RFO(\mathbf{q},\mathbf{q}',\Dz)
&:=\Rzero(\mathbf{q},\mathbf{q}',\Dz,0)\,
\EStot(\mathbf{q},\mathbf{q}',\Dz)\,
\ECtot(\mathbf{q},\mathbf{q}'),
\label{eq:RFO}\\
\RCTF(\mathbf{q},\mathbf{q}',\Dz)
&:=\Rzero(\mathbf{q},\mathbf{q}',\Dz,0)\,
\EStot(\mathbf{q},\mathbf{0},\Dz)\,
\EStot(\mathbf{q}',\mathbf{0},\Dz)\,
\ECtot(\mathbf{q},\mathbf{0})\,
\ECtot(\mathbf{q}',\mathbf{0}),
\label{eq:RCTF}
\end{align}
where the last two chromatic factors are \emph{not} conjugated, matching
printed~(7a).
Spatial envelopes are real, so conjugating $\EStot(\mathbf{q}',\mathbf{0},\Dz)$
is redundant. The spatial CTF factors still depend on defocus.
\end{definition}

\begin{lemma}[Hermiticity of $\Rzero$]
\label{lem:R0herm}
$\conj{\Rzero(\mathbf{q},\mathbf{q}',\Dz,\varepsilon)}
=\Rzero(\mathbf{q}',\mathbf{q},\Dz,\varepsilon)$.
\end{lemma}
\begin{proof}
$M$ is real. Each of $W_S$ and $W_C$ is unimodular, so
$W^*(\mathbf{q})W(\mathbf{q}')=\conj{W(\mathbf{q})W^*(\mathbf{q}')}$
for $W\in\{W_S,W_C\}$. The product of these factors therefore swaps
$(\mathbf{q},\mathbf{q}')$ under conjugation.
\end{proof}

\begin{theorem}[FO recovers CTF on the axis]
\label{thm:axis-eq}
$\RFO(\mathbf{q},\mathbf{0},\Dz)=\RCTF(\mathbf{q},\mathbf{0},\Dz)$.
\end{theorem}
\begin{proof}
On $\mathbf{q}'=\mathbf{0}$, \cref{thm:axes} gives $\Gamma_C=\Gamma_S=1$,
or expand~\eqref{eq:RFO}--\eqref{eq:RCTF} using
$\EStot(\mathbf{0},\mathbf{0},\Dz)=1$ and $\ECtot(\mathbf{0},\mathbf{0})=1$.
\end{proof}

\begin{theorem}[Paper comparison identity]
\label{thm:R-gamma}
\begin{equation}
\label{eq:R-gamma}
\boxed{\RCTF=\RFO\,\Gamma_C\,\Gamma_S.}
\end{equation}
\end{theorem}
\begin{proof}
Neither envelope vanishes:
$\EStot$ is an exponential, and $\ECtot=\ECC$ times an exponential with
$\ECC=z^{-1/2}$ and $z\neq 0$ ($\ReOp(z)=1$).
Multiplying~\eqref{eq:RFO} by~\eqref{eq:paper7} therefore cancels
$\EStot(\mathbf{q},\mathbf{q}')$ and $\ECtot(\mathbf{q},\mathbf{q}')$,
leaving~\eqref{eq:RCTF}.
\end{proof}

Thus $\Gamma_C\Gamma_S\neq 1$ is precisely the warning that a separable CTF
kernel differs from the FO kernel.

\section{Jacobi--Anger expansion of a sinusoidal object}
\label{sec:jacobi}

A kinematic LEEM phase object of physical amplitude $A$ and wavelength
$\Lambda$ is
$\psi_0(x)\propto\exp(\ii\varphi\sin(kx))$ with
$\varphi=4\pi A/\lambda_0$ and $k=2\pi/\Lambda$.
The phase amplitude is a definition of the reflection kinematics, not an
FO theorem.

\begin{definition}[Fourier--Bessel coefficients]
\label{def:bessel}
For $\varphi\in\R$ and $n\in\Z$,
\begin{equation}
\label{eq:Jn}
J_n(\varphi)
:=\frac{1}{2\pi}
\int_{(0,2\pi]}
\exp(\ii\varphi\sin\theta)\,\ee^{-\ii n\theta}\,\dd\theta.
\end{equation}
\end{definition}

\begin{lemma}[Periodicity and unit modulus]
\label{lem:phase}
$\theta\mapsto\exp(\ii\varphi\sin\theta)$ is $2\pi$-periodic and has
modulus $1$.
Its second derivative is bounded by $\abs{\varphi}+\varphi^2$.
\end{lemma}
\begin{proof}
Sine is $2\pi$-periodic.
The exponent $\ii\varphi\sin\theta$ is purely imaginary, so the exponential
has modulus $1$.
Differentiating,
$f'=\ii\varphi\cos\theta\cdot f$ and
$f''=\ii\varphi(-\sin\theta)f+(\ii\varphi\cos\theta)^2 f$,
hence $\abs{f''}\le\abs{\varphi}+\varphi^2$.
\end{proof}

\begin{lemma}[Bessel bound]
\label{lem:bessel-bound}
If $n\neq 0$ then $\abs{J_n(\varphi)}\le(\abs{\varphi}+\varphi^2)/n^2$.
Consequently $\sum_{n\in\Z}J_n(\varphi)$ is absolutely summable.
\end{lemma}
\begin{proof}
Two integrations by parts on the circle (periodic endpoints cancel)
express $J_n$, for $n\neq 0$, as the Fourier coefficient of $f''$
divided by $(in)^2$.
The bound on $\abs{f''}$ yields the coefficient bound.
Comparison with $\sum n^{-2}$ gives summability; the $n=0$ term is finite.
\end{proof}

\begin{theorem}[Jacobi--Anger]
\label{thm:jacobi}
For all real $\varphi,\theta$,
\begin{equation}
\label{eq:jacobi}
\exp(\ii\varphi\sin\theta)
=\sum_{n\in\Z}J_n(\varphi)\,\ee^{\ii n\theta},
\end{equation}
the series converging absolutely and uniformly on $\R$.
Hence $\exp(\ii\varphi\sin(kx))=\sum_n J_n(\varphi)\,\ee^{\ii nkx}$.
\end{theorem}
\begin{proof}
The map $\theta\mapsto\exp(\ii\varphi\sin\theta)$ is continuous and
$2\pi$-periodic.
Its Fourier coefficients~\eqref{eq:Jn} are absolutely summable
by~\cref{lem:bessel-bound}.
A continuous function on the circle with absolutely summable Fourier
coefficients equals its Fourier series pointwise on $(0,2\pi]$.
Both sides of~\eqref{eq:jacobi} are $2\pi$-periodic, so the identity
extends to all real $\theta$ by reducing modulo $2\pi$
($\ee^{\ii n(\theta-2\pi m)}=\ee^{\ii n\theta}$).
\end{proof}

\begin{corollary}[Discrete spectrum]
\label{cor:spectrum}
In the paper's $2\pi$ Fourier convention,
$\ee^{\ii n k x}=\ee^{2\pi\ii q_n x}$ with $q_n=nk/(2\pi)=n/\Lambda$.
The object is a discrete spectrum at those frequencies.
Larger $\varphi$ populates larger $\abs{n}$, so pairs $(q,q')$ leave the
axes where $\Gamma_C\Gamma_S=1$, and a separable CTF becomes inaccurate.
\end{corollary}

\section{Assembly}
\label{sec:assembly}

\begin{theorem}[Working formula]
\label{thm:assembly}
Under~\cref{approx:taylor},
\begin{equation}
\label{eq:assembly}
R(\mathbf{q},\mathbf{q}',\Dz)
=\Rzero(\mathbf{q},\mathbf{q}',\Dz,0)\,
\EStot(\mathbf{q},\mathbf{q}',\Dz)\,
\ECtot(\mathbf{q},\mathbf{q}'),
\end{equation}
with $\EStot$ from~\cref{thm:ES-fwhm} and $\ECtot$ from~\cref{thm:EC-closed}.
Nac/ac specialisation is substitution of coefficients, not a new theorem.
\end{theorem}
\begin{proof}
Substitute~\eqref{eq:R0approx} into~\eqref{eq:R} and use that the
$(\mathbf{k},\varepsilon)$ integral factors.
\end{proof}

\section{Correspondence with the Lean formalization}
\label{sec:lean}

The identities above are encoded without $\mathtt{sorry}$ in this
repository, using Mathlib's Gaussian and complex-power lemmas.
Wavelength is the field \texttt{lam} (Lean reserves $\lambda$).

\begin{center}
\small
\begin{tabular}{@{}lll@{}}
\toprule
Statement & Paper & Lean name \\
\midrule
FWHM / $\sigma^2=\Delta^2/(8\ln 2)$ & --- & \texttt{gaussian\_fwhm}, \texttt{variance\_from\_fwhm} \\
Char.\ fun.\ $E[\ee^{2\pi\ii a k}]$ & A1 & \texttt{charFun\_source1D}, \texttt{charFun\_source2D} \\
$\nabla\chi_S$ / $\partial_q\chi_S$ & --- & \texttt{hasDerivAt\_chiS}, \texttt{hasGradientAt\_chiS2} \\
Taylor remainder $o(k)$ & (3) & \texttt{chiS\_taylor} \\
$b_1\varepsilon+b_2\varepsilon^2$ & --- & \texttt{chiC\_sub} \\
$E_S=\ee^{-\kappa a^2}$ & (5) & \texttt{spatialEnvelope\_eq\_fwhm} \\
$\ECtot$ closed form & (6a)--(6b) & \texttt{chromaticEnvelopeIntegral\_eq\_closed} \\
Polar $\ECC$ & A2 & \texttt{ecc\_eq\_polar} \\
$\Gamma_S=\ee^{-2\kappa uv}$ & (7b) & \texttt{gammaS\_eq\_kappa}, \texttt{gammaS\_ac} \\
Axes / perfect coherence & --- & \texttt{gammaS\_axis}, \texttt{gamma\_product\_perfect} \\
$\RCTF=\RFO\Gamma_C\Gamma_S$ & (7) & \texttt{R\_CTF\_eq\_R\_FO\_mul\_gamma} \\
Jacobi--Anger & \S3.2 & \texttt{jacobi\_anger} \\
\bottomrule
\end{tabular}
\end{center}

Modelling approximations~\ref{approx:aperture} and~\ref{approx:taylor} are
\emph{not} Lean theorems: the working kernel $\RFO$ is defined as the
factorised expression after those truncations.

\section{Analytic pitfalls}
\label{sec:pitfalls}

\begin{enumerate}[leftmargin=*,itemsep=0.35em]
\item \textbf{Principal square root.}
Always $\exp(\tfrac12\mathrm{Log}\,z)$ with $\Arg\in(-\pi,\pi]$.
The identity $(A/z)^{1/2}=A^{1/2}z^{-1/2}$ requires $\Arg(z)\neq\pi$.
For $z=1-4\pi\ii b_2\sE^2$ one has $\ReOp(z)=1$, so the cut is avoided
and $b_2=0$ yields $+1$.
\item \textbf{Argument of $1-\ii y$.}
$\Arg(1-\ii y)=-\arctan y$ for \emph{all} real $y$ because the real part
is identically $1$.
There is no branch jump as $y$ crosses $0$.
\item \textbf{Complex quadratic Gaussians.}
\Cref{lem:cgaussian} requires $\ReOp(B)<0$, which comes only from the
real Gaussian density.
A purely imaginary quadratic coefficient would make the integral
oscillatory.
\item \textbf{Sign of $-C^2/(4B)$.}
$C=2\pi\ii b_1$ implies $C^2=-4\pi^2 b_1^2$.
Forgetting $\ii^2=-1$ produces a purely oscillatory ``envelope''.
\item \textbf{Two-dimensional inner product.}
$\abs{\mathbf{a}}^2$ contains $\mathbf{q}\cdot\mathbf{q}'$, not $q q'$.
$\Gamma_S$ contains $\nabla\chi_S(\mathbf{q})\cdot\nabla\chi_S(\mathbf{q}')$,
not a product of Euclidean norms.
\item \textbf{Printed $\Gamma_C$ has no extra conjugate.}
On the axis $\mathbf{q}=\mathbf{0}$, only $\abs{\Gamma_C}=1$ in general
for ac.
\item \textbf{FWHM versus variance.}
$\kappa=\pi^2\qill^2/(4\ln 2)$ equals $2\pi^2\sill^2$, not $\pi^2\sigma^2$.
Mixing FWHM and $\sigma$ typically produces a factor-of-four error.
\item \textbf{Defocus sign.}
$W=\ee^{+2\pi\ii\chi}$ with $+\tfrac12\Dz\lambda q^2$.
Do not convert $\chi$ to radians and then insert another $2\pi$.
\end{enumerate}

\begin{thebibliography}{9}
\bibitem{Yu2019}
K.~M. Yu, K.~L.~W. Lau, M.~S. Altman,
\emph{Fourier optics of image formation in aberration-corrected LEEM},
Ultramicroscopy \textbf{200} (2019) 160--168.

\bibitem{Pang2009}
A.~B. Pang, Th.~M\"uller, M.~S. Altman, E.~Bauer,
J.~Phys.: Condens.\ Matter \textbf{21} (2009) 314006.

\bibitem{Schramm2012}
S.~M. Schramm, A.~B. Pang, M.~S. Altman, R.~M. Tromp,
Ultramicroscopy \textbf{115} (2012) 88--108.

\bibitem{Yu2017}
K.~M. Yu, A.~Locatelli, M.~S. Altman,
Ultramicroscopy \textbf{183} (2017) 109--116.

\bibitem{Tromp2012}
R.~M. Tromp, S.~M. Schramm,
Ultramicroscopy (2012), optimization of the CTF.

\bibitem{Hanen1971}
K.~J. Han{\ss}en, L.~Trepte, Optik \textbf{32} (1971) 519.
\end{thebibliography}

\end{document}
